% Chapter 4, Section 1 _Linear Algebra_ Jim Hefferon
%  http://joshua.smcvt.edu/linearalgebra
%  2001-Jun-12
\chapter{行列式} 
在第一章中，我们
重点讨论了方程个数与未知量个数相同的线性方程组
这一特殊情形，即形如 \( T\vec{x}=\vec{b} \)、其中 $T$ 为方阵的方程组。
我们注意到，这样的 $T$ 只有两种。
若 $T$ 对任意 $\vec{b}$，
都对应唯一解，比如
齐次方程组 $T\vec{x}=\zero$ 那样，那么
对每一个这样的 $\vec{b}$，$T$ 都对应唯一解。
我们称这样的矩阵为非奇异的。
另一类 $T$，凡是以其为
系数矩阵的线性方程组都要么无解、
要么有无穷多解，
我们称之为奇异的。

在此后的工作中，这一区分一直是一条主线。
例如，我们现在知道，一个
\( \nbyn{n} \) 矩阵 \( T \) 是非奇异的当且仅当
以下每一条都成立：
%<*EquivalentOfNonsingular>
\begin{itemize}
   \item 任意方程组 \( T\vec{x}=\vec{b} \) 都有解，
          且该解唯一；
    \item 对 $T$ 作高斯–若尔当消元得到单位矩阵；
    \item $T$ 的各行构成一个线性无关的集合；
    \item \( T \) 的各列构成一个线性无关的集合，
          即 \( \Re^n \) 的一个基；
    \item \( T \) 所表示的任意映射都是同构；
    \item 存在逆矩阵 \( T^{-1} \)。
\end{itemize}
%</EquivalentOfNonsingular>
因此，当我们看到一个方阵时，首先
要问的问题之一就是它是否
非奇异。

本章将建立一个公式，用来判定 $T$ 是否非奇异。
更准确地说，我们将对 $\nbyn{1}$~矩阵
给出一个公式，对 $\nbyn{2}$~矩阵给出一个，等等。
这些公式自然相互关联；也就是说，
我们将建立一族
公式，一个描述每种规格之公式的方案。

由于我们的讨论仅限于方阵，本章中
我们常常径直说“矩阵”，以代替“方阵”。





\section{定义}
%<*DeterminantIntro>
判定非奇异性这件事，
对 \( \nbyn{1} \) 矩阵而言是平凡的。
\begin{equation*}
  \begin{mat}
       a
  \end{mat} 
  \quad\text{是非奇异的当且仅当}\quad
  a \neq 0
\end{equation*} 
推论~Three.IV.\ref{cor:TwoByTwoInv} 给出了 $\nbyn{2}$ 的公式。 
\begin{equation*}
     \begin{mat}
           a  &b  \\
           c  &d
      \end{mat}  
   \quad\text{是非奇异的当且仅当}\quad
   ad-bc \neq 0
\end{equation*}
我们可以像前面那样得出 $\nbyn{3}$ 的公式，
尽管计算相当繁琐
% \typeout{DET1 ThreeByThreeDetForm cleveref expmt!}
% (see \cref{exer:ThreeByThreeDetForm}).
（见\nearbyexercise{exer:ThreeByThreeDetForm}）。
\begin{equation*}
     \begin{mat}
           a  &b  &c  \\
           d  &e  &f  \\
           g  &h  &i
     \end{mat}
  \quad\text{是非奇异的当且仅当}\quad
    aei+bfg+cdh-hfa-idb-gec \neq 0
\end{equation*}
心怀这些特例，我们设想一族
公式：$a$，$ad-bc$，等等。
对每个 $n$，该公式定义了一个 
\definend{行列式}\index{determinant}\index{matrix!determinant} 
函数
$\map{\det_{\nbyn{n}}}{\matspace_{\nbyn{n}}}{\Re}$ 
使得 $\nbyn{n}$ 矩阵 $T$ 是非奇异的当
且仅当 $\det_{\nbyn{n}}(T)\neq 0$。
%</DeterminantIntro>
（我们通常省略下标 \( \nbyn{n} \)，因为
\( T \) 的规格已经表明我们指的是哪个行列式函数。）






\subsectionoptional{探索}
\textit{本小节为选读内容，
给出一般定义的动机与展开。
定义在下一小节中。}

在前文中，各种情形下的矩阵都是非奇异的当且
仅当某个公式不为零。
但这三个公式并
没有显示出明显的规律。
我们或许会注意到 \(\nbyn{1}\) 的项
\( a \) 有一个字母，\(\nbyn{2}\) 的项
\(ad\) 与 \(bc\) 各有两个字母，而 \(\nbyn{3}\)
的项各有三个字母。
我们甚至能注意到，在这些项中
每一项都是从矩阵的每一行与每一列各取一个字母，例如，
在 \(cdh\) 项中一个字母
来自每一行、每一列。
\begin{equation*}
   \begin{mat}
          &    &c \\
      d           \\
          &h
   \end{mat} 
\end{equation*}
但这些观察与其说是有启发性
的，不如说是令人困惑的。
比如，我们可能想知道为什么
有些项相加，而有些项相减。

解题的一个好策略是
先探索解必须具有哪些性质，然后
再去寻找具有这些性质的东西。
于是我们先来问：我们希望行列式
公式具有哪些性质。

就目前而言，
判定一个矩阵是否奇异的主要方法
是作高斯
消元，然后检查
阶梯形矩阵的对角线上是否有零，
即沿对角线的乘积是否为~零。
于是我们可以猜想：无论找到怎样的行列式公式，证明
它正确的做法可能都涉及对矩阵应用高斯消元法，
以说明最终沿对角线的乘积为零当且仅当
我们的公式给出零。

这提示了一个计划：~我们将寻找一族行列式
公式，它们
不受行变换影响，且阶梯形
矩阵的行列式是其对角元之积。
% Under this plan, a proof that the functions determine singularity would go, 
% ``Where $T\rightarrow\cdots\rightarrow\hat{T}$ is the Gaussian
% reduction, the determinant of $T$ equals the
% determinant of $\hat{T}$ (because the determinant is unchanged by row
% operations), which is the product down the diagonal, which is
% zero if and only if the matrix is singular''.
在本小节余下的部分，我们将用 $\nbyn{2}$ 与
$\nbyn{3}$ 的公式来检验这一计划。
最后我们将不得不修改“不受行变换影响”
这一部分，但改动不大。

首先我们检查
$\nbyn{2}$ 与 $\nbyn{3}$ 的公式是否不受
倍加这一行变换的影响：~若
\begin{equation*}
   T \grstep{k\rho_i+\rho_j} \hat{T}
\end{equation*}
那么 \( \det(\hat{T})=\det(T) \) 是否成立？ 
对 $k\rho_1+\rho_2$ 操作之后的 $\nbyn{2}$ 行列式作这一检验
\begin{equation*}
   \det(
       \begin{mat}
         a     &b       \\
         ka+c  &kb+d    \\
      \end{mat}
   )
   = a(kb+d)-(ka+c)b = ad-bc
\end{equation*}
表明它确实不变，而
另一个 $\nbyn{2}$ 倍加 $k\rho_2+\rho_1$ 给出同样的结果。
同样，
$\nbyn{3}$ 的倍加 $k\rho_3+\rho_2$ 也使行列式保持不变
  \begin{align*}
    \det(
      \begin{mat}
         a    &b    &c    \\
         kg+d &kh+e &ki+f \\
         g    &h    &i
      \end{mat}
    )
    &=\begin{array}[t]{@{}l@{}}
         a(kh+e)i+b(ki+f)g+c(kg+d)h \\
         \ \hbox{}-h(ki+f)a-i(kg+d)b-g(kh+e)c  
    \end{array}                                 \\
     &=aei + bfg + cdh - hfa - idb - gec
  \end{align*}
其他的 $\nbyn{3}$ 行倍加操作也是如此。

所以这个计划看起来是有希望的。
当然，也许如果我们算出
$\nbyn{4}$ 行列式公式并加以检验，可能会发现
它会受行倍加的影响。
这只是一次探索，我们还没有掌握全部事实。
尽管如此，到目前为止一切顺利。

接下来，对于行对换的情形，我们比较 \( \det(\hat{T}) \) 与
\( \det(T) \)。
% \begin{equation*}
%    T \grstep{ {\rho}_i \leftrightarrow {\rho}_j } \hat{T}
% \end{equation*}
这里我们遇到了一个障碍：
\(\nbyn{2}\) 的行对换 $\rho_1\leftrightarrow\rho_2$
并不给出 \( ad-bc \)。
  \begin{equation*}
     \det(
       \begin{mat}
         c  &d \\
         a  &b
       \end{mat}
     )
     = bc - ad
  \end{equation*}
而 \(\nbyn{3}\) 矩阵内的这次 $\rho_1\leftrightarrow\rho_3$ 对换
\begin{equation*}
   \det(
     \begin{mat}
        g  &h  &i \\
        d  &e  &f \\
        a  &b  &c
     \end{mat}
   )
   = gec + hfa + idb - bfg - cdh - aei
\end{equation*}
同样不给出与对换前相同的行列式，因为又一次
出现了符号改变。
再试 \(\nbyn{3}\) 的另一次对换 $\rho_1\leftrightarrow\rho_2$ 
\begin{equation*}
   \det(
     \begin{mat}
        d  &e  &f \\
        a  &b  &c \\
        g  &h  &i
     \end{mat}
   )
   = dbi + ecg + fah - hcd - iae - gbf
\end{equation*}
也给出符号的改变。

所以在这一实验中，行对换
看来会改变行列式的符号。
这并没有完全破坏我们的计划。
我们希望通过只考虑
公式是否给出零来判定非奇异性，而不考虑其符号。
因此，与其期望行列式公式
完全不受行变换影响，不如修改计划：在行对换时
它们改变符号。

显然，
我们最后要比较 \( \det(\hat{T}) \) 与 \( \det(T) \) 
针对的操作
% \begin{equation*}
%    T \grstep{ k{\rho}_i } \hat{T}
% \end{equation*}
是把一行乘以一个标量。 
下面这个计算
\begin{equation*}
  \det(
    \begin{mat}
      a   &b   \\
      kc  &kd
    \end{mat}
  )
  = a(kd) - (kc)b
  =k\cdot (ad-bc)
\end{equation*}
最后整个行列式乘上了~$k$，
另一个 $\nbyn{2}$ 情形也得到同样的结果。
这个 \(\nbyn{3}\) 情形以同样的方式收尾
\begin{align*}
  \det(
  \begin{mat}
    a    &b    &c   \\
    d    &e    &f   \\
    kg   &kh   &ki
  \end{mat}
  )
   &= \begin{array}[t]{@{}l@{}}
         ae(ki) + bf(kg) + cd(kh)                \\
         \>- (kh)fa - (ki)db - (kg)ec  
      \end{array}                                      \\
   &= k\cdot(aei + bfg + cdh - hfa - idb - gec)
\end{align*}
另外两个 $\nbyn{3}$ 情形也是如此。
这些使我们猜想：把一行乘以~$k$
会把行列式乘以~$k$。
与前面一样，这修改了我们的计划，但没有破坏它。
我们只要求
行列式公式的为零性保持不变，而不关注
它的符号或大小。

所以在这次探索中，我们的计划
经过一些非实质性的修改，现在变为：~我们将寻找
$\nbyn{n}$ 行列式
函数，它们在行倍加操作下
保持不变，在行对换时
改变符号，在某行被倍乘时按比例缩放，
并且阶梯形矩阵的行列式
是沿对角线的乘积。
在下面两小节中我们将看到，对每个~$n$，
这样的函数有且仅有一个。

最后，为下一小节作准备，请注意提出标量因子是一种
逐行进行的操作：
在这里
\begin{equation*}
    \det(
      \begin{mat}[r]
        3  &3  &9  \\
        2  &1  &1  \\
        5  &11 &-5
     \end{mat}
    )
    =3 \cdot \det(
               \begin{mat}[r]
                 1  &1  &3  \\
                 2  &1  &1  \\
                 5  &11 &-5
               \end{mat}
             )                                  
\end{equation*}
$3$ 仅从最上面一行提出，其余各行保持不变。
因此，在行列式的定义中，我们将
把它写成各行的函数
\( \det (\vec{\rho}_1,\vec{\rho}_2,\dots\vec{\rho}_n) \)，而不是
\( \det(T) \)，也不是各元素
\( \det(t_{1,1},\dots,t_{n,n}) \) 的函数。

